

\section{Tree-shaking}


\subsection{Syntactic Tree shaking}

\begin{enumerate}
	\item This describes optimizations applied to the automated TLA+ code to remove any unused TLA+ code
	\item There are two main sources of unused code: Unused boilerplate macros for operations like transitive closure
	\item Unused macros involving user imports, which are all defined and translated, but never used in Init
	\item Both these classes of code are removed via post-processing on the TLA+ code
	\item The expected increase is a small improvement in the time taken by the SANY parser check, as well as optimizations in the Z3 translation output from Apalache
	\item These optimizations, if turned off, have no effect on the output of the translated code
\end{enumerate}

\subsection{Semantic Tree shaking}

\begin{enumerate}
	\item These optimizations involve the removal of unused elements from the base Alloy model
	\item This requires knowledge about the underlying Alloy model being translated, and is not solely an operation done on TLA+
	\item Since Apalache and TLC do not understand Alloy or the mappings between the Alloy code and the translator output, these optimizations can be performed only by the translator itself
	\item Unused fields and sigs are removed entirely from the TLA+ model, by purging the associated variables and by removing the clauses that use these variables in the memebrship constraints on Init, and by modifying Next to avoid referring to these variables
	\item This may uncover possible further syntactic tree-shaking, so this is performed first. Syntactic tree-shaking cannot reveal further semantic tree-shaking, so the processes are not cycled till the code converges into a fixpoint.
	\item The application of this optimization changes the output produced by the translated model, by eliminating references to certain variables
\end{enumerate}

